\documentclass[11pt]{article}
\usepackage[margin=1in]{geometry}
\usepackage{amsmath,amssymb,amsthm,mathtools}
\usepackage{booktabs,graphicx,xcolor,hyperref}
\usepackage{microtype}
\usepackage[numbers,sort&compress]{natbib}

\newtheorem{theorem}{Theorem}
\newtheorem{proposition}[theorem]{Proposition}
\newtheorem{corollary}[theorem]{Corollary}
\newtheorem{lemma}[theorem]{Lemma}
\theoremstyle{definition}
\newtheorem{definition}[theorem]{Definition}
\theoremstyle{remark}
\newtheorem{remark}[theorem]{Remark}

\newcommand{\Pzero}{\mathbb P_0}
\newcommand{\Ezero}{\mathbb E_0}
\newcommand{\TESS}{\operatorname{TESS}}
\newcommand{\ind}{\mathbbm{1}}

\title{Inferential Search Size of Adaptive Statistical Policies}
\author{Atsushi Senda}
\date{}

\begin{document}
\maketitle

\begin{abstract}
% 150--200 words. Lead with non-identifiability, then sharp policy calculus,
% threshold coherence/reversal, two-bank inference, and the locked illustration.
\end{abstract}

\section{Introduction}
% Adaptive candidate sets; missing policy layer; relation to fixed-family
% effective numbers, data snooping, post-selection and valid adaptive testing.
% State contributions in one paragraph.

\section{Policy-level inferential search size}
\subsection{Definition and characterization}
For a declared policy $\mathcal A$ under global-null law $P_0$, define
\[
\pi_{\mathcal A}(\alpha)=\Pzero\{P_{\mathcal A}<\alpha\},\qquad
S_{\mathcal A}(\alpha)=
\frac{\log\{1-\pi_{\mathcal A}(\alpha)\}}{\log(1-\alpha)}.
\]

\begin{proposition}[Independent-composition characterization]
% State null-law invariance, independent-union additivity, continuity,
% exact-single normalization.
\end{proposition}

\subsection{Fixed-family bridge}
\begin{remark}[Copula diagonal and extremal coefficient]
% Keep short; explicitly state prior-art boundary.
\end{remark}

\section{Adaptive-policy representation and sharp envelopes}
Let $R_0(\alpha)$ and $R_1(\alpha)$ denote counterfactual base and
expanded rejection indicators, $D_\alpha=R_1(\alpha)-R_0(\alpha)$,
and let $S$ contain activation-time information. For activation kernel
$a(S)\in[0,1]$, define $m_\alpha(S)=\Ezero(D_\alpha\mid S)$.

\begin{theorem}[Policy representation and sharp same-budget envelope]
% pi_a = pi_0 + E[a m]; upper/lower integrated quantiles at fixed rate r.
\end{theorem}

\begin{corollary}[Identifiability]
% Same candidate law and activation rate identify rejection burden iff m is a.s. constant.
\end{corollary}

\begin{corollary}[Expected-budget insufficiency]
% Fixed expected candidate count does not upper-bound deep-tail search size.
\end{corollary}

\section{Threshold coherence and policy reversal}
\begin{definition}[Rank coherence]
% Formal definition with separability/common-version conditions.
\end{definition}

\begin{theorem}[Common-rank simultaneous optimality]
% One rank ordering is optimal over alpha interval and all rates.
\end{theorem}

\begin{theorem}[Same-law, same-budget curve reversal]
% State either for a general branch-output policy or provide a nested-expansion construction.
\end{theorem}

\section{Two-bank inference}
\subsection{Population and finite-reference targets}
% Distinguish population-calibrated, conditional finite-reference, and
% unconditional two-bank sampling targets.

\subsection{Estimated-trigger expansion}
\begin{theorem}[Estimated-trigger two-bank asymptotic linearity]
% Reference IF: a0 W0 + b1 W1 + dU WU; evaluation IF; asymptotic normality.
\end{theorem}

\begin{theorem}[Complete-replication two-bank bootstrap]
% State bootstrap consistency and TESS delta-method result.
\end{theorem}

\subsection{Paired budget-matched policy contrast}
Let $e_0=E(R_0)$, $r=E(A)$, $\mu=E(D)$ and $\nu=E(AD)$. Then
\[
\pi_A=e_0+\nu,\qquad
\pi_C=e_0+r\mu,\qquad
\nu-r\mu=\operatorname{Cov}(A,D).
\]

\begin{theorem}[Paired budget-matched TESS contrast]
% D4 target: joint evaluation/reference IF and paired bootstrap validity.
\end{theorem}

\subsection{Finite-reference and plus-one bridge}
\begin{lemma}[Strict plus-one rejection as an order-statistic event]
% Exact integer rounding, no-tie form, and first-order equivalence regime.
\end{lemma}

\section{Locked machine-learning illustration}
% One compact design paragraph; primary contrast; rescue reversal; signed
% mechanism; two-bank sensitivity. SUPPORT2 is brief here and detailed in supplement.

\section{Discussion}
% Coupling property, expected-budget impossibility, threshold non-orderability,
% diagnostic vs error control, regular model vs exact pipeline, reporting implications.

\section*{Supplementary material}
% Full proofs, D1/D2 validation, cost/FOSD/weighted extensions, empirical details,
% SUPPORT2, manifests, software and archive provenance.

\bibliographystyle{plainnat}
\bibliography{references}
\end{document}
